% =========================================================================
% English translation (generated by AI using Claude Opus 4.8 (Max effort)).
% This file was translated from Hungarian to English by AI.
% DISCLAIMER: Please review against the original Hungarian .tex files, before relying on it!
% SOURCE: 13.1 Logika.tex
% Note: in the truth table the Hungarian "i"/"h" (igaz/hamis) were rendered
%       as "T"/"F" (true/false). CNF/DNF abbreviations: KNF->CNF, DNF->DNF,
%       KDNF->canonical DNF (CDNF), KKNF->canonical CNF (CCNF).
% =========================================================================
\documentclass[margin=0px]{article}

\usepackage{listings}
\usepackage[utf8]{inputenc}
\usepackage{graphicx}
\usepackage{float}
\usepackage[a4paper, margin=0.7in]{geometry}
\usepackage{amsthm}
\usepackage{amssymb}
\usepackage{fancyhdr}
\usepackage{setspace}

\onehalfspacing

\makeatletter
\renewcommand\paragraph{%
	\@startsection{paragraph}{4}{0mm}%
	{-\baselineskip}%
	{.5\baselineskip}%
	{\normalfont\normalsize\bfseries}}
\makeatother

\renewcommand{\figurename}{Figure}

\newenvironment{tetel}[1]{\paragraph{#1 \\}}{}

\pagestyle{fancy}
\lhead{\it{CS BSc Final Exam Topics}}
\rhead{13.1 Logic}

\title{\textbf{{\Large ELTE IK - Computer Science BSc} \vspace{0.2cm} \\ {\huge Final Exam Topics}} \vspace{0.3cm} \\ 13.1 Logic}
\author{}
\date{}

\begin{document}
\maketitle

\begin{center}
\fbox{\parbox{0.92\textwidth}{\small \textit{Note: This document was translated from Hungarian to English by AI. Please verify the information against the original Hungarian source before relying on it.}}}
\end{center}

\begin{tetel}{Logic}
    Propositional calculus and first-order predicate calculus: syntax, semantics, equivalent transformations, the concept of semantic consequence, resolution.
\end{tetel}

\section{Logic}

\subsection{Basic concepts}

The subject of logic is the examination of the human thinking process and
the search for, or creation of, correct forms of thinking.\\

\noindent Concepts:
\begin{enumerate}
    \item	\textbf{Statement}: A declarative sentence whose logical value (truth) is
          decidable, true or false in any context. We say
          that a statement is true if its information content corresponds to reality (the facts),
          and false in the opposite case.

          The declarative sentences used in everyday speech are most often not statements, since
          the context is also part of the meaning of the sentence: time, state of the environment, a certain level
          of general knowledge, etc. (e.g. ``the sun was shining at 8 a.m. this morning'' is not a statement, but e.g.
          ``every even number is divisible by 2'' is a statement).

    \item	\textbf{Truth value}: The set of truth values is $\mathbb{L}=\left\{true,false\right\}$.

    \item	\textbf{Form of thinking}: By a form of thinking we mean a pair $(F,A)$ where
          $A$ is a statement, and $F=\left\{A_{1},A_{2},...,A_{n}\right\}$ is a set of statements.

          The form of thinking is correct if in every case when every statement of $F$ is
          true, $A$ is also true.
\end{enumerate}

\subsection{Propositional calculus}

\subsubsection{The syntax of propositional logic}

\paragraph{The alphabet of propositional logic}
The alphabet of propositional logic is $V_{0}=V_{v} \cup \left\{(,)\right\} \cup \left\{\neg, \wedge, \vee, \supset\right\}$, where $V_{v}$ is the
set of propositional variables. So $V_{0}$ contains the propositional variables, the parentheses, and the signs of the logical operations.

\paragraph{The language of propositional logic}
The language of propositional logic ($\mathcal{L}_{0}$) consists of propositional-logic formulas, which are produced as follows:
\begin{enumerate}
    \item	Every propositional variable is a propositional-logic formula. These are the so-called prime formulas (or atomic formulas).

    \item	If $A$ is a propositional-logic formula, then so is $\neg A$.

    \item	If $A$ and $B$ are propositional-logic formulas, then $(A \wedge B)$, $(A \vee B)$ and $(A \supset B)$ are also
          propositional-logic formulas.

    \item	Every propositional-logic formula is produced by the finitely many times application of rules 1-3.
\end{enumerate}

\noindent \textbf{Literal}: If $X$ is a propositional variable, then the formulas $X$ and $\neg X$ are literals whose base is $X$.\\

\noindent \textbf{Direct subformula}:
\begin{enumerate}
    \item	A prime formula has no direct subformula.

    \item	The direct subformula of $\neg A$ is $A$.

    \item	The direct subformulas of $A \circ B$ ($\circ$ is one of the binary connectives $\wedge, \vee, \supset$) are $A$ (left-hand)
          and $B$ (right-hand).
\end{enumerate}

\noindent \textbf{Subformula}: Let $A \in \mathcal{L}_{0}$ be a propositional-logic formula. Then the set of subformulas of $A$
is the narrowest set that
\begin{enumerate}
    \item	has $A$ as an element, and

    \item	if the formula C is an element, then the direct subformulas of C are also elements.

\end{enumerate}

\noindent \textbf{Structure tree}: The structure tree of a formula C is a finite ordered tree whose vertices are formulas,
\begin{enumerate}
    \item	its root is C,

    \item	the vertex $\neg A$ has exactly one child, $A$,

    \item	the vertex $A \circ B$ has exactly two children, respectively the formulas $A$ and $B$,

    \item	its leaves are prime formulas.
\end{enumerate}

\begin{figure}[H]
    \centering
    \includegraphics[width=0.3\linewidth]{../../13. Számításelmélet (és Logika)/img/szerkfa}
    \caption{Example of a structure tree.}
    \label{fig:szerkfa}
\end{figure}

\noindent \textbf{Logical complexity}: The logical complexity of a formula is the number of logical connectives found in it.\\

\noindent \textbf{Scope of an operation}: The scope of an operation is, among the subformulas of the formula, the
subformula of smallest logical complexity in which the given operation occurs.\\

\noindent \textbf{Main logical connective}: The main logical connective of a formula is the connective whose
scope is the formula itself.\\

\noindent \textbf{Precedence}: The precedence of the logical connectives in decreasing order is the following: $\neg, \wedge, \vee, \supset$.\\

Based on the definitions, it is unambiguous in a \textit{fully parenthesized formula} what the scope of the logical connectives is and what the main
logical connective is. Now we show in which cases and which parentheses bounding which subformulas can be omitted in a formula so that the scope of the logical connectives does not change. Among the subformulas, the prime formulas and the negation formulas have no outer parenthesis pair, so we only have to decide about subformulas of the form $(A \circ B)$ whether $A \circ B$ can be written instead. We always start the omission of parentheses with the omission of the outer parenthesis pair of the formula (if there is one). Then, once we have examined the question of outer-parenthesis omission in a subformula, after that we deal with the outer parentheses of the direct subformulas of this subformula. Two cases are possible:

\begin{enumerate}
    \item	The subformula is a negation formula, in which the outer parentheses of the direct subformula of the form $(A \circ B)$ cannot be omitted.

    \item	The subformula is a formula of the form $(A \bullet B)$ or $A \bullet B$, in whose direct subformulas $A$ and $B$ the fate of the outer parentheses has to be decided. If the formula $A$ is of the form $A_{1} \circ A_{2}$, then the outer parenthesis pair of $A$ can be omitted if $\circ$ has higher precedence than $\bullet$. If the formula $B$ is of the form $B_{1} \circ B_{2}$, then the outer parenthesis pair of $B$ can be omitted if $\circ$ has higher or equal precedence than $\bullet$.

    \item	If some direct subformula of a formula of the form $(A \wedge B)$ or $A \wedge B$ is likewise a conjunction, or some direct subformula of a formula of the form
          $(A \vee B)$ or $A \vee B$ is likewise a disjunction, then the outer parenthesis pair can be omitted from such subformulas.
\end{enumerate}

\noindent \textbf{Formula chains}: Taking into account the conventions on the omission of parentheses, we can also obtain so-called conjunction, disjunction, and implication formula chains. Their form is $A_{1} \wedge ... \wedge A_{n}$, $A_{1} \vee ... \vee A_{n}$, or
$A_{1} \supset ... \supset A_{n}$. The main logical connective of these chain formulas we will determine with the following parenthesization convention: $(A_{1} \wedge (A_{2} \wedge ... \wedge (A_{n-1} \wedge A_{n})...))$, $(A_{1} \vee (A_{2} \vee ... \vee (A_{n-1} \vee A_{n})...))$, or $(A_{1} \supset (A_{2} \supset ... \supset (A_{n-1} \supset A_{n})...))$

\subsubsection{The semantics of propositional logic}

\noindent \textbf{Interpretation}: By an interpretation of $\mathcal{L}_{0}$ we mean a function $\mathcal{I} : V_{v} \to \mathbb{L}$ that unambiguously assigns a truth value to every propositional variable.\\

\noindent \textbf{Boolean valuation}: The \textit{Boolean valuation in interpretation $\mathcal{I}$} of formulas in $\mathcal{L}_{0}$ is the following function $\mathcal{B}_{\mathcal{I}} : \mathcal{L}_{0} \to \mathbb{L}$:

\begin{enumerate}
    \item if $A$ is a prime formula, then $\mathcal{B}_{\mathcal{I}}(A) = \mathcal{I}(A)$,

    \item let $\mathcal{B}_{\mathcal{I}}(\neg A)$ be $\neg \mathcal{B}_{\mathcal{I}}(A) $,

    \item let $\mathcal{B}_{\mathcal{I}}(A \wedge B)$ be $\mathcal{B}_{\mathcal{I}}(A) \wedge \mathcal{B}_{\mathcal{I}}(B)$,

    \item let $\mathcal{B}_{\mathcal{I}}(A \vee B)$ be $\mathcal{B}_{\mathcal{I}}(A) \vee \mathcal{B}_{\mathcal{I}}(B)$,

    \item let $\mathcal{B}_{\mathcal{I}}(A \supset B)$ be $\mathcal{B}_{\mathcal{I}}(A) \supset \mathcal{B}_{\mathcal{I}}(B)$,
\end{enumerate}

\noindent \textbf{Basis}: A fixed order of the propositional variables of the formula.\\

\noindent \textbf{Semantic tree}: The different interpretations of a formula can be illustrated with the help of a semantic tree. The semantic
tree is a binary tree whose $i$-th level ($i>=1$) belongs to the $i$-th propositional variable of the basis, and from every vertex two edges start, one
labeled with the propositional variable assigned to the level, the other with its negation. For the propositional variable $X$, let the label $X$ mean that $X$ is \textit{true} in the given interpretation, and the label $\neg X$ that it is \textit{false} in the given interpretation. Every branch of the semantic tree
represents a possible interpretation. For a formula of $n$ variables every branch is $n$ long, and the tree has $2^{n}$ branches and contains all possible interpretations.

\begin{figure}[H]
    \centering
    \includegraphics[width=0.6\linewidth]{../../13. Számításelmélet (és Logika)/img/szemantikusfa}
    \caption{The semantic tree of the formula containing the propositional variables X, Y, Z.}
    \label{fig:szemantikusfa}
\end{figure}

\noindent \textbf{Truth table}: The truth table of a formula of $n$ variables is a table consisting of $n+1$ columns and $2^{n}$ rows.
In the header of the table, to the $i$-th column ($1<=i<=n$) the $i$-th propositional variable of the formula's basis, and to the $(n+1)$-th column the formula itself
is assigned. In the first $n$ columns, for the individual rows, we give respectively the different interpretations of the formula, then in the column of the formula
we write into every row the truth value of the formula -- obtained by the Boolean valuation in the interpretation belonging to the row.\\

\noindent \textbf{The truth table of the logical operations}:

\begin{table}[H]
    \begin{tabular}{ll|lllll}
        X & Y & $\neg X$ & $X \wedge Y$ & $X \vee Y$ & $ X \supset Y$ & \\ \hline
        T & T & F        & T            & T          & T              & \\
        T & F & F        & F            & T          & F              & \\
        F & T & T        & F            & T          & T              & \\
        F & F & T        & F            & F          & T              &
    \end{tabular}
\end{table}

\noindent \textbf{True set, false set}: The true set $(A^{i})$ of a formula $A$
is the set of those interpretations on which the truth valuation of the formula is true. And the false set $(A^{h})$ of the formula $A$
is the set of those interpretations for which the truth valuation of the formula is false.\\

\noindent \textbf{Truth valuation function}: A function that assigns to every formula its true set ($\varphi A^{i}$) or
its false set ($\varphi A^{h}$).\\

Let $A$ be an arbitrary propositional-logic formula. Let us determine for $A$ the $\varphi A^{i}$ and
$\varphi A^{h}$ conditions relating to its interpretations as follows:

\begin{enumerate}
    \item	If $A$ is a prime formula, the condition $\varphi A^{i}$ is satisfied exactly by those interpretations $\mathcal{I}$ for which
          $\mathcal{I}(A)=true$, and the condition $\varphi A^{h}$ exactly by those for which $\mathcal{I}(A)=false$.

    \item	The conditions $\varphi (\neg A)^{i}$ hold exactly when the conditions $\varphi A^{h}$ hold.

    \item	The conditions $\varphi (A \wedge B)^{i}$ hold exactly when the conditions $\varphi A^{i}$ and $\varphi B^{i}$ hold simultaneously.

    \item	The conditions $\varphi (A \vee B)^{i}$ hold exactly when the conditions $\varphi A^{i}$ or $\varphi B^{i}$ hold.

    \item	The conditions $\varphi (A \supset B)^{i}$ hold exactly when the conditions $\varphi A^{h}$ or $\varphi B^{i}$ hold.
\end{enumerate}

\noindent \textbf{Theorem}: For any propositional-logic formula $A$, the conditions $\varphi A^{i}$ are satisfied exactly by the interpretations in $A^{i}$.\\

\noindent \textbf{Truth-valuation tree}: The interpretations of a formula $A$ satisfying the $\varphi A^{i}$ or $\varphi A^{h}$ conditions
can be illustrated with the help of the truth-valuation tree. We produce the truth-valuation tree using the structure tree of the formula.
To the root we assign which truth-valuation conditions of $A$ for which truth value we are looking for, then under the root the direct subformulas of $A$ get placed with the appropriate condition prescription, as follows:

\begin{figure}[H]
    \centering
    \includegraphics[width=0.7\linewidth]{../../13. Számításelmélet (és Logika)/img/igazsagertfa}
    \caption{Condition prescriptions of the truth-valuation tree.}
    \label{fig:igazsagertfa}
\end{figure}

After this we assign to the root the \checkmark (processed) mark. We continue the procedure recursively until, on a branch, the unprocessed
formulas

\begin{itemize}
    \item	(a) all become propositional variables, or

    \item	(b) a mutually contradicting prescription appears for the same formula.
\end{itemize}

In case (a) the $n$-tuples containing the truth values, fixed on the branch, of the propositional variables occurring on the branch are all elements of the true set of the formula in the case of a $\varphi A^{i}$ root, and of the false set of the formula in the case of a $\varphi A^{h}$ root.

In case (b) no such truth-value $n$-tuple is produced.

\begin{figure}[H]
    \centering
    \includegraphics[width=0.5\linewidth]{../../13. Számításelmélet (és Logika)/img/igazsagertfapelda}
    \caption{The truth-valuation tree of the formula $(Y \vee Z) \wedge (Z \supset \neg X)$.}
    \label{fig:igazsagertfapelda}
\end{figure}

\noindent In the above example the true set of the formula, based on the truth-valuation tree, is: $\left\{(T,T,F),(F,T,T),(F,T,F),(F,F,T)\right\}$\\

\noindent \textbf{Extended truth table}: To facilitate the computation of the truth value of the formula in a truth table, the extended truth table was introduced. In the extended truth table, besides the columns assigned to the propositional variables and to the formula, the columns belonging to the subformulas of the formula also appear, in turn. Essentially the subformulas appearing in the structure tree are listed.

\begin{figure}[H]
    \centering
    \includegraphics[width=0.5\linewidth]{../../13. Számításelmélet (és Logika)/img/kiterjigaztabla}
    \caption{The extended truth table of the formula $(Y \vee Z) \wedge (Z \supset \neg X)$.}
    \label{fig:kiterjigaztabla}
\end{figure}

\noindent \textbf{Satisfiability of a formula, its model}: A propositional-logic formula $A$ is \textit{satisfiable} if there exists an interpretation $\mathcal{I}$
for which $\mathcal{I} \models_{0} A$, that is, the Boolean valuation $\mathcal{B}_{\mathcal{I}}$ assigns a true value to $A$. Such
an interpretation is called a \textit{model} of $A$. If $A$ has no model, then we say that it is \textit{unsatisfiable}.

If in the truth table of $A$ there is a row in which a true value appears in the column of the formula, then the formula is satisfiable, otherwise unsatisfiable. Likewise, if $\varphi A^{i}$ is not empty, then it is satisfiable, otherwise unsatisfiable.\\

\noindent \textbf{Propositional-logic law, tautology}: A propositional-logic formula $A$ is a \textit{propositional-logic law} or otherwise a \textit{tautology} if every interpretation of $\mathcal{L}_{0}$ is a model of $A$. (notation: $\models_{0} A$)\\

\noindent \textbf{Decision problem}: We call the following tasks decision problems:
\begin{enumerate}
    \item	Decide for an arbitrary formula whether it is a tautology!

    \item	Decide for an arbitrary formula whether it is unsatisfiable!
\end{enumerate}

\noindent \textbf{Tautologically equivalent formulas}: The propositional-logic formulas $A$ and $B$ are \textit{tautologically equivalent} (notation: $A \sim _{0} B$) if in every interpretation $\mathcal{I}$ of $\mathcal{L}_{0}$, $\mathcal{B}_{\mathcal{I}}(A)=\mathcal{B}_{\mathcal{I}}(B)$.\\

\noindent \textbf{Satisfiability of a formula set, its model}: An arbitrary set $\Gamma$ of formulas of $\mathcal{L}_{0}$ is
satisfiable if there is an interpretation $\mathcal{I}$ of $\mathcal{L}_{0}$ for which: $\forall A \in \Gamma: \mathcal{I} \models_{0} A$.
Such an interpretation $\mathcal{I}$ is a model of $\Gamma$. If $\Gamma$ has no model, then $\Gamma$ is unsatisfiable.\\

\noindent \textbf{Lemma}: The models of a formula set $\left\{A_{1},A_{2},...,A_{n}\right\}$ are exactly those interpretations $\mathcal{I}$
that are models of the formula $A_{1} \wedge A_{2} \wedge ... \wedge A_{n}$. Consequently $\left\{A_{1},A_{2},...,A_{n}\right\}$
is unsatisfiable exactly when the formula $A_{1} \wedge A_{2} \wedge ... \wedge A_{n}$ is unsatisfiable.\\

\noindent \textbf{Semantic consequence}: Let $\Gamma$ be an arbitrary set of propositional-logic formulas, $B$ an arbitrary formula.
We say that the formula $B$ is a \textit{tautological consequence} of the formula set $\Gamma$ (notation: $\Gamma \models_{0} B$) if every interpretation that is a model of $\Gamma$ is also a model of $B$. The formulas in $\Gamma$ are called condition formulas or premises,
and the formula B the consequence formula (conclusion).\\

\noindent \textbf{Theorem}: Let $\Gamma$ be an arbitrary set of propositional-logic formulas, $A$, $B$, $C$ arbitrary propositional-logic formulas.
If $\Gamma \models_{0} A$, $\Gamma \models_{0} B$ and $\left\{A,B\right\} \models_{0} C$, then $\Gamma \models_{0} C$.\\

\noindent \textbf{Theorem}: Let $A_{1},A_{2},...,A_{n}, B$ be arbitrary propositional-logic formulas.
$\left\{A_{1},A_{2},...,A_{n}\right\} \models_{0} B$ exactly when the formula set $\left\{A_{1},A_{2},...,A_{n}, \neg B\right\}$
is unsatisfiable, that is, the formula $A_{1} \wedge A_{2} \wedge ... \wedge A_{n} \wedge \neg B$ is unsatisfiable.\\

\noindent \textbf{Theorem}: Let $A_{1},A_{2},...,A_{n}, B$ be arbitrary propositional-logic formulas.
$\left\{A_{1},A_{2},...,A_{n}\right\} \models_{0} B$ exactly when
$\models_{0} A_{1} \wedge A_{2} \wedge ... \wedge A_{n} \supset B$.\\

\paragraph{Equivalent transformations}
Concepts:
\begin{enumerate}
    \item	A prime formula (propositional variable), or its negation, is by a common name called a \textit{literal}. The prime formula
          is the \textit{base of the literal}. A literal is in certain cases also called a \textit{unit conjunction} or \textit{unit disjunction}
          (\textit{unit clause}).

    \item	An \textit{elementary conjunction} is the unit conjunction, or the conjunction of literals with different bases ($\wedge$ connection
          between the literals). An \textit{elementary disjunction} or \textit{clause} is the unit disjunction and the disjunction of literals with different bases
          ($\vee$ connection between the literals). An elementary conjunction, or elementary disjunction, is \textit{complete}
          with respect to an $n$-variable logical operation if all $n$ propositional variables are the base of some of its literals.

    \item	By \textit{disjunctive normal form} (DNF) we mean the disjunction of elementary conjunctions.
          By \textit{conjunctive normal form} (CNF) we mean the conjunction of elementary disjunctions. We speak of \textit{canonical}
          disjunctive or conjunctive normal forms (CDNF, or CCNF) if the
          elementary conjunctions, or elementary disjunctions, appearing in them are complete.

\end{enumerate}

\noindent \textbf{Producing the CDNF, CCNF describing an arbitrary logical operation}: Let $b: \mathbb{L}^{n} \to \mathbb{L}$
be an $n$-variable logical operation. Let us give the operation table of $b$.
In the header of the first $n$ columns let us write the propositional variables $X_{1}$, $X_{2}$, ... $X_{n}$.\\

\noindent Producing the CDNF describing $b$:

\begin{enumerate}
    \item	Let us choose those rows in the operation table where $b$ assigns a \textit{true}
          value to the given truth-value $n$-tuple. Let these rows be respectively $s_{1}, s_{2}, .. s_{r}$. To every such row let us assign
          a complete elementary conjunction $X_{1}' \wedge X_{2}' \wedge ... \wedge X_{n}'$ such that the literal $X_{j}'$
          is $X_{j}$ or $\neg X_{j}$ according to whether in this row $X_{j}$ has truth value \textit{true} or \textit{false}.
          Let the complete elementary conjunctions thus obtained be respectively $k_{s_{1}}, k_{s_{2}}, .. k_{s_{r}}$.

    \item	From the complete elementary conjunctions thus obtained let us make a disjunction chain formula:
          $k_{s_{1}} \vee k_{s_{2}} \vee ... \vee k_{s_{r}}$. This formula will be the canonical disjunctive
          normal form (CDNF) of the operation $b$.
\end{enumerate}

\begin{figure}[H]
    \centering
    \includegraphics[width=0.5\linewidth]{../../13. Számításelmélet (és Logika)/img/kdnf_pelda}
    \caption{The operation table of a three-variable logical operation $b$ and the produced complete elementary conjunctions.}
    \label{fig:kdnf_pelda}
\end{figure}

The canonical disjunctive normal form of the operation $b$ in the above example is the following formula:\\
$(X \wedge Y \wedge \neg Z) \vee (X \wedge \neg Y \wedge \neg Z) \vee (\neg X \wedge Y \wedge Z) \vee (\neg X \wedge \neg Y \wedge Z) \vee (\neg X \wedge \neg Y \wedge \neg Z)$.\\

\noindent Producing the CCNF describing $b$:

\begin{enumerate}
    \item	Let us choose those rows in the operation table where $b$ assigns a \textit{false}
          value to the given truth-value $n$-tuple. Let these rows be respectively $s_{1}, s_{2}, .. s_{r}$. To every such row let us assign
          a complete elementary disjunction $X_{1}' \vee X_{2}' \vee ... \vee X_{n}'$ such that the literal $X_{j}'$
          is $X_{j}$ or $\neg X_{j}$ according to whether in this row $X_{j}$ has truth value \textit{false} or \textit{true}.
          Let the complete elementary disjunctions thus obtained be respectively $d_{s_{1}}, d_{s_{2}}, .. d_{s_{r}}$.

    \item	From the complete elementary disjunctions thus obtained let us make a conjunction chain formula:
          $d_{s_{1}} \wedge d_{s_{2}} \wedge ... \wedge d_{s_{r}}$. This formula will be the canonical conjunctive
          normal form (CCNF) of the operation $b$.
\end{enumerate}

\begin{figure}[H]
    \centering
    \includegraphics[width=0.5\linewidth]{../../13. Számításelmélet (és Logika)/img/kknf_pelda}
    \caption{The operation table of a three-variable logical operation $b$ and the produced complete elementary disjunctions.}
    \label{fig:kknf_pelda}
\end{figure}

The canonical conjunctive normal form of the operation $b$ in the above example is the following formula:\\
$(\neg X \vee \neg Y \vee \neg Z) \wedge (\neg X \vee Y \vee \neg Z) \wedge (X \vee Y \vee \neg Z)$. \\

\noindent \textbf{Simplification of CNF, DNF}: By the logical complexity of a propositional-logic formula we meant the number of
logical connectives appearing in the formula. Among formulas describing the same logical operation, we regard as simpler the one
whose logical complexity is smaller (that is, which contains fewer logical connectives).

Let $X$ be a propositional variable, $k$ an elementary conjunction not containing $X$, $d$ an elementary
disjunction not containing $X$. Then, applying the simplification rules

\begin{itemize}
    \item	(a) $(X \wedge k) \vee (\neg X \wedge k) \sim_{0} k $ and

    \item	(b) $(X \vee d) \wedge (\neg X \vee d) \sim_{0} d $
\end{itemize}

we can rewrite conjunctive and disjunctive normal forms into a simpler form.\\

\noindent Classical Quine--McCluskey algorithm for simplifying CDNF:

\begin{enumerate}
    \item	Let us list all the complete elementary conjunctions appearing in the CDNF in the list $L_{0}$, $j:=0$.

    \item	We examine all the possible elementary-conjunction pairs appearing in $L_{j}$, whether the
          simplification rule (a) is applicable to them. If yes, then we mark the two chosen conjunctions with $\checkmark$,
          and write the resulting conjunction into the list $L_{j+1}$. Those elementary conjunctions that during the examination of $L_{j}$
          do not get marked were not simplifiable, so they get into the simplified disjunctive
          normal form.

    \item	If the conjunction list $L_{j+1}$ is not empty, then $j:=j+1$. Let us perform step 2 again.

    \item	From the elementary conjunctions obtained during the algorithm but not marked, let us make a disjunction
          chain formula. Thus we obtain a simplified DNF logically equivalent to the original CDNF.
\end{enumerate}

\paragraph{Resolution}

Let $A_{1}, A_{2}, ... , A_{n}, B$ be arbitrary propositional-logic formulas.
We want to prove that $\left\{A_{1}, A_{2}, ... , A_{n}\right\} \models_{0} B$,
which is equivalent to $\left\{A_{1}, A_{2}, ... , A_{n}, \neg B\right\}$ being unsatisfiable.
Let us rewrite the formulas of this latter formula set into CNF form! Then we obtain the
formula set $\left\{CNF_{A_{1}}, CNF_{A_{2}}, ... , CNF_{A_{n}}, CNF_{\neg B}\right\}$,
which is unsatisfiable exactly when the set of clauses appearing in the formulas of the set is
unsatisfiable.

According to the simplification rule for clauses, if $X$ is a propositional variable, and $C$ is a clause not
containing $X$, then $(X \vee C) \wedge (\neg X \vee C) \sim_{0} C$. The clause equivalent to the conjunction of the
unit clauses $X$ and $\neg X$ (we say that $X$ and $\neg X$ are a complementary literal pair) that is
simpler and contains not a single literal is the empty clause, which we denote by the $\square$ sign
and which by definition has the truth value false in every interpretation.

Let now $C_{1}$ and $C_{2}$ be clauses that contain exactly one complementary literal pair,
that is, $C_{1} = C_{1}' \vee L_{1}$ and $C_{2} = C_{2}' \vee L_{2}$, where $L_{1}$ and $L_{2}$ are the single
complementary literal pair ($C_{1}'$ and $C_{2}'$ can also be empty clauses). It is clear that if in the two clauses
there are also literals besides the complementary literal pair, and these are not all identical, the applicability condition of the simplification
rule does not hold.\\

\noindent \textbf{Theorem}: If $C_{1} = C_{1}' \vee L_{1}$ and $C_{2} = C_{2}' \vee L_{2}$,
where $L_{1}$ and $L_{2}$ are a complementary literal pair, then $\left\{C_{1},C_{2}\right\} \models_{0} C_{1}' \vee C_{2}'$\\

\noindent \textbf{Resolvent}: Let $C_{1}$ and $C_{2}$ be clauses that contain exactly one complementary literal pair,
that is, $C_{1} = C_{1}' \vee L_{1}$ and $C_{2} = C_{2}' \vee L_{2}$, where $L_{1}$ and $L_{2}$ are the complementary literal pair;
the clause $C_{1}' \vee C_{2}'$ is called the \textit{resolvent} of the clause pair $(C_{1},C_{2})$ (or of the formula $C_{1} \vee C_{2}$).
If $C_{1} = L_{1}$ and $C_{2} = L_{2}$ (that is, $C_{1}'$ and $C_{2}'$ are empty clauses), their resolvent is the empty clause ($\square$).
The activity whose result is the resolvent is \textit{resolving}.

\begin{figure}[H]
    \centering
    \includegraphics[width=0.7\linewidth]{../../13. Számításelmélet (és Logika)/img/rezolvens}
    \caption{Examples of the resolvability and resolvent of clause pairs.}
    \label{fig:rezolvens}
\end{figure}

\noindent \textbf{Theorem}: If the clause $C$ is the resolvent of the clause pair $(C_{1},C_{2})$, then those interpretations $\mathcal{I}$
cannot satisfy the clause set $\left\{C_{1}, C_{2}\right\}$ in which the truth value of $C$ is false, that is,
$\mathcal{B}_{\mathcal{I}}(C) = false$.\\

\noindent \textbf{Resolution derivation}: A resolution derivation of the clause $C$ from a clause set $S$ is a finite
clause sequence $k_{1}, k_{2}, ... ,k_{m} (m \geq 1)$ where for every $j = 1, 2, ..., m$

\begin{enumerate}
    \item	either $k_{j} \in S$,

    \item	or there is some $1 \leq s,t \le j$ such that $k_{j}$ is the resolvent of the clause pair $(k_{s}, k_{t})$,
\end{enumerate}

\noindent and the last term of the clause sequence, $k_{m}$, is exactly the clause $C$.

By our convention the decision problem of the resolution calculus is whether $\square$ is derivable from $S$. The
goal of the resolution derivation is therefore the derivation of $\square$ from $S$. That $\square$ is derivable from $S$
can also be expressed as: there exists a resolution refutation of $S$.\\

\noindent Example: Let us try to derive $\square$ from the clause set
$S = \left\{\neg X \vee Y, \neg Y \vee Z, X \vee Z, \neg V \vee Y \vee Z, \neg Z \right\}$.
The derivation can be started from any clause in $S$.

\begin{figure}[H]
    \centering
    \includegraphics[width=0.4\linewidth]{../../13. Számításelmélet (és Logika)/img/rezoluciopelda}
    \caption{Resolution derivation of $\square$ from $S$.}
    \label{fig:rezoluciopelda}
\end{figure}

\noindent \textbf{Lemma}: Let $S$ be an arbitrary clause set. In the case of a resolution derivation from $S$, any clause derived from $S$
is a tautological consequence of $S$.\\

\noindent \textbf{Soundness of the resolution calculus}: The resolution calculus is \textit{sound}, that is, for any
clause set $S$, if $\square$ is derivable from $S$, then $S$ is \textit{unsatisfiable}.\\

\noindent \textbf{Completeness of the resolution calculus}: The resolution calculus is \textit{complete}, that is, for any finite, unsatisfiable
clause set $S$, $\square$ is derivable from $S$.\\

\noindent \textbf{Derivation tree}: The structure of a resolution derivation can be illustrated with the help of a \textit{derivation tree}.
The vertices of the derivation tree are clauses. From two vertices an edge leads into a third, common vertex exactly when that is the resolvent of the two clauses.

\begin{figure}[H]
    \centering
    \includegraphics[width=0.5\linewidth]{../../13. Számításelmélet (és Logika)/img/rezoluciopelda_levfa}
    \caption{The derivation tree of the previous example.}
    \label{fig:rezoluciopelda_levfa}
\end{figure}

\noindent \textbf{Resolution strategies}:

\begin{itemize}
    \item	\textit{Linear resolution}: A linear resolution derivation from a clause set $S$ is a
          resolution derivation $k_{1},l_{1},k_{2},l_{2}, ..., k_{m-1},l_{m-1}, k_{m}$ in which for every
          $j = 2, 3, ..., m$, $k_{j}$ is the resolvent of the clause pair $(k_{j-1},l_{j-1})$. The clauses $k_{j}$ we call central
          clauses, and the clauses $l_{j}$ side clauses.

          Any resolution derivation can be rewritten into a linear one, that is, the linear resolution calculus is complete.

    \item	\textit{Linear input resolution}: A linear input resolution derivation from a clause set $S$
          is a linear resolution derivation $k_{1},l_{1},k_{2},l_{2}, ..., k_{m-1},l_{m-1}, k_{m}$ in which for every
          $j = 1, 2, ..., m-1$, $l_{j} \in S$, that is, in the linear input resolution derivation the side clauses are elements of $S$.

          The linear input resolution strategy is not complete, but a formula class can be given for which it is. The clauses containing at most one
          negated literal are called Horn clauses, and the Horn formulas are those formulas whose
          conjunctive normal form is a conjunction of Horn clauses. The linear input resolution strategy is complete in the case of Horn formulas.
\end{itemize}

\subsection{Predicate calculus}

\subsubsection{Syntax of first-order logical languages}

The alphabet of a first-order logical language contains logical and non-logical symbols, as well as separators.
The non-logical symbol set can be given in the form $<Srt, Pr, Fn, Cnst>$, where:

\begin{enumerate}
    \item	$Srt$ is a nonempty set, its elements symbolize sorts,

    \item	$Pr$ is a nonempty set, its elements are predicate symbols,

    \item	the elements of the set $Fn$ are function symbols,

    \item	and $Cnst$ is the set of constant symbols.
\end{enumerate}

\noindent The signature of the alphabet $<Srt, Pr, Fn, Cnst>$ is a triple $<\nu_{1}, \nu_{2}, \nu_{3}>$, where

\begin{enumerate}
    \item	to every predicate symbol $P \in Pr$, $\nu_{1}$ assigns the shape of the predicate symbol,
          that is, the sort sequence $(\pi_{1}, \pi_{2}, ..., \pi_{k})$,

    \item	to every function symbol $f \in Fn$, $\nu_{2}$ assigns the shape of the function symbol,
          that is, the sort sequence $(\pi_{1}, \pi_{2}, ..., \pi_{k}, \pi)$ and

    \item	to every constant symbol $c \in Cnst$, $\nu_{3}$ assigns the shape of the constant symbol,
          that is, $(\pi)$
\end{enumerate}

\noindent ($k > 0$ and $\pi_{1}, \pi_{2}, ..., \pi_{k}, \pi \in Srt$).

Logical signs are the logical connectives also used in propositional logic, as well as the universal ($\forall$)
and existential ($\exists$) quantifiers and the individual variables of different sorts. In a first-order
language alphabet, to every sort $\pi \in Srt$ belongs a countably infinite system
$v_{1}^{\pi}, v_{2}^{\pi}, ...$ of symbols; these symbols we call variables of sort $\pi$.
Separators are the opening and closing parentheses, and the comma.

In first-order logical languages the separators and the logical signs are always the same, but
the set of non-logical signs, and their signature, can differ significantly from language to language.
Therefore we always give the quadruple $<Srt, Pr, Fn, Cnst>$ and its $<\nu_{1}, \nu_{2}, \nu_{3}>$
signature when we refer to the alphabet of a first-order logical language. Its notation is $V[V_{\nu}]$,
where $V_{\nu}$ gives the quadruple $<Srt, Pr, Fn, Cnst>$ with signature $<\nu_{1}, \nu_{2}, \nu_{3}>$.\\

\noindent \textbf{Terms}: The set of terms over the alphabet $V[V_{\nu}]$ is $\mathcal{L}_{t}[V_{\nu}]$, which
has the following properties:

\begin{enumerate}
    \item	Every variable and constant of sort $\pi \in Srt$ is a term of sort $\pi$.

    \item	If the function symbol $f \in Fn$ is of the form $(\pi_{1}, \pi_{2}, ..., \pi_{k}, \pi)$
          and $t_{1}, t_{2}, ..., t_{k}$ are terms -- respectively of sorts $\pi_{1}, \pi_{2}, ..., \pi_{k}$ --,
          then $f(s_{1}, s_{2}, ..., s_{k})$ is a term of sort $\pi$.

    \item	Every term is produced by the finitely many times application of rules 1-2.
\end{enumerate}

\noindent \textbf{Formulas}: The set of first-order formulas over the alphabet $V[V_{\nu}]$ is $\mathcal{L}_{f}[V_{\nu}]$, which
has the following properties:

\begin{enumerate}
    \item	If the predicate symbol $P \in Pr$ is of the form $(\pi_{1}, \pi_{2}, ..., \pi_{k})$
          and $t_{1}, t_{2}, ..., t_{k}$ are terms -- respectively of sorts $\pi_{1}, \pi_{2}, ..., \pi_{k}$ --,
          then the word $P(t_{1}, t_{2}, ..., t_{k})$ is a first-order formula.
          The formulas thus obtained are called atomic formulas.

    \item	If $S$ is a first-order formula, then so is $\neg S$.

    \item	If $S$ and $T$ are first-order formulas and $\circ$ is a binary logical connective,
          then $(S \circ T)$ is also a first-order formula.

    \item	If $S$ is a first-order formula, $Q$ is a quantifier ($\forall$ or $\exists$) and $x$
          is an arbitrary variable, then $QxS$ is also a first-order formula. The formulas thus obtained are called quantified formulas,
          the formulas of the form $\forall xS$ are universally quantified formulas, and the formulas of the form $\exists xS$
          are existentially quantified formulas. In quantified formulas $Qx$ is the prefix of the formula, and $S$
          its matrix.

    \item	Every first-order formula is produced by the finitely many times application of rules 1-4.
\end{enumerate}

\noindent The first-order logical language over the alphabet $V[V_{\nu}]$ is
$\mathcal{L}[V_{\nu}] = \mathcal{L}_{t}[V_{\nu}] \cup \mathcal{L}_{f}[V_{\nu}]$, that is, every word of $\mathcal{L}[V_{\nu}]$
is either a term or a formula.\\

The negation, conjunction, disjunction, implication (their meaning is the same as in propositional logic) and quantified formulas
are compound formulas.

The prime formulas of the first-order logical language are the atomic formulas and the quantified formulas.\\

\noindent \textbf{Kinds of variable occurrence}: An occurrence of variable $x$ of a formula is:

\begin{itemize}
    \item	free, if it does not fall within the scope of a quantifier relating to $x$,

    \item	bound, if it falls within the scope of a quantifier relating to $x$.
\end{itemize}

\noindent \textbf{Kinds of variable}: A variable $x$ of a formula is:

\begin{itemize}
    \item	free, if all of its occurrences are free,

    \item	bound, if all of its occurrences are bound, and

    \item	mixed, if it has both free and bound occurrences.
\end{itemize}

\noindent \textbf{Closedness, openness of a formula}: A formula is:

\begin{itemize}
    \item	closed, if all of its variables are bound,

    \item	open, if at least one of its variables has a free occurrence and

    \item	quantifier-free, if there is no quantifier in it
\end{itemize}

Remark: closed formulas symbolize first-order statements (a first-order statement is nothing other than a declarative sentence formulated for a set of elements).

\subsubsection{The semantics of first-order logic}

\noindent \textbf{Mathematical structure}: By a mathematical structure we mean a quadruple $<U, R, M, K>$, where:

\begin{enumerate}
    \item	$U = \bigcup_{\pi} U_{\pi}$ is a nonempty base set (universe),

    \item	$R$ is the set of logical functions (relations) defined on $U$,

    \item	$M$ is the set of mathematical functions (basic operations) defined on $U$,

    \item	$K$ is the set of distinguished elements (constants) of $U$ (can be empty).
\end{enumerate}

\noindent \textbf{Interpretation}: The interpretation is a mathematical structure $<U, R, M, K>$ and
a function quadruple $\mathcal{I} = <\mathcal{I}_{Srt}, \mathcal{I}_{Pr}, \mathcal{I}_{Fn}, \mathcal{I}_{Cnst}>$,
where:

\begin{itemize}
    \item	the function $\mathcal{I}_{Srt} : \pi \mapsto U_{\pi}$ gives to every sort $\pi \in Srt$
          a nonempty set $U_{\pi}$, the set of individuals of sort $\pi$,

    \item	the function $\mathcal{I}_{Pr} : P \mapsto P^{\mathcal{I}}$ gives to every
          predicate symbol $P \in Pr$ of the form $(\pi_{1}, \pi_{2}, ..., \pi_{k})$ a
          logical function (relation)
          $P^{\mathcal{I}} : U_{\pi_{1}} \times U_{\pi_{2}} \times ... \times U_{\pi_{k}} \to \mathbb{L}$,

    \item	the function $\mathcal{I}_{Fn} : f \mapsto f^{\mathcal{I}}$ assigns to every
          function symbol $f \in Fn$ of the form $(\pi_{1}, \pi_{2}, ..., \pi_{k}, \pi)$ a
          mathematical function (operation)
          $P^{\mathcal{I}} : U_{\pi_{1}} \times U_{\pi_{2}} \times ... \times U_{\pi_{k}} \to U_{\pi}$,

    \item	and $\mathcal{I}_{Cnst} : c \mapsto ct^{\mathcal{I}}$ assigns to every constant symbol
          $c \in Cnst$ of sort $\pi$ an individual of the individual domain $U_{\pi}$,
          that is, $c^{\mathcal{I}} \in U_{\pi}$.
\end{itemize}

\noindent \textbf{Variable evaluation}: Let $\mathcal{I}$ be an interpretation of the language $\mathcal{L}[V_{\nu}]$,
let the universe of the interpretation be U and let V denote the set of variables of the language. A mapping $\kappa : V \to U$,
where if $x$ is a variable of sort $\pi$, then $\kappa(x) \in U_{\pi}$, is called a variable evaluation in $\mathcal{I}$.\\

\noindent \textbf{Semantics of $\mathcal{L}_{t}[V_{\nu}]$}:
Let $\mathcal{I}$ be an interpretation of the language $\mathcal{L}[V_{\nu}]$ and $\kappa$ a variable evaluation in $\mathcal{I}$.
The value of a term $t$ of sort $\pi$ of the language $\mathcal{L}[V_{\nu}]$ in $\mathcal{I}$
under the variable evaluation $\kappa$ is the following individual in $U_{\pi}$ -- denoted by $|t|^{\mathcal{I},\kappa}$:

\begin{enumerate}
    \item	if $c \in Cnst$ is a constant symbol of sort $\pi$, then $|c|^{\mathcal{I},\kappa}$ is the individual $c^{\mathcal{I}}$ in $U_{\pi}$,

    \item	if $x$ is a variable of sort $\pi$, then $|x|^{\mathcal{I},\kappa}$ is the individual $\kappa(x)$ in $U_{\pi}$,

    \item	if $t_{1}, t_{2}, ..., t_{k}$ are terms respectively of sorts $\pi_{1}, \pi_{2}, ..., \pi_{k}$ and their values
          under the variable evaluation $\kappa$ are respectively the individuals $|t_{1}|^{\mathcal{I},\kappa}$ in $U_{\pi_{1}}$, the
          $|t_{2}|^{\mathcal{I},\kappa}$ in $U_{\pi_{2}}$ ... and $|t_{k}|^{\mathcal{I},\kappa}$ in $U_{\pi_{k}}$,
          then for a function symbol $f \in Fn$ of the form $(\pi_{1}, \pi_{2}, ..., \pi_{k}, \pi)$,
          $|f(t_{1}, t_{2}, ..., t_{k})|^{\mathcal{I},\kappa}$ is the individual
          $f^{\mathcal{I}}(|t_{1}|^{\mathcal{I},\kappa}, |t_{2}|^{\mathcal{I},\kappa}, ..., |t_{k}|^{\mathcal{I},\kappa})$ in $U_{\pi}$.

\end{enumerate}

\noindent \textbf{$x$-variant of a variable evaluation}: Let $x$ be a variable. The variable evaluation $\kappa^{*}$ is an
$x$-variant of the variable evaluation $\kappa$ if $\kappa^{*}(y) = y$ for every variable $y$ different from $x$.\\

\noindent \textbf{Logical value of a first-order logical formula}:
Let $\mathcal{I}$ be an interpretation of the language $\mathcal{L}[V_{\nu}]$ and $\kappa$ a variable evaluation in $\mathcal{I}$.
To a formula $C$ of the language $\mathcal{L}[V_{\nu}]$ we assign, in $\mathcal{I}$ under the variable evaluation $\kappa$,
the following truth value -- denoted by $|C|^{\mathcal{I},\kappa}$:

\begin{enumerate}
    \item	$|P(t_{1}, t_{2}, ..., t_{k})|^{\mathcal{I},\kappa} = \left\{
              \begin{array}{lr}
                  true  & : P^{\mathcal{I}}(|t_{1}|^{\mathcal{I},\kappa}, |t_{2}|^{\mathcal{I},\kappa}, ..., |t_{k}|^{\mathcal{I},\kappa}) = true \\
                  false & : \text{otherwise}
              \end{array}
              \right\}$

    \item	let $|\neg A|^{\mathcal{I},\kappa}$ be $\neg |A|^{\mathcal{I},\kappa}$

    \item	let $|A \wedge B|^{\mathcal{I},\kappa}$ be $|A|^{\mathcal{I},\kappa} \wedge |B|^{\mathcal{I},\kappa}$

    \item	let $|A \vee B|^{\mathcal{I},\kappa}$ be $|A|^{\mathcal{I},\kappa} \vee |B|^{\mathcal{I},\kappa}$

    \item	let $|A \supset B|^{\mathcal{I},\kappa}$ be $|A|^{\mathcal{I},\kappa} \supset |B|^{\mathcal{I},\kappa}$

    \item	$|\forall x A|^{\mathcal{I},\kappa} = \left\{
              \begin{array}{lr}
                  true  & : |A|^{\mathcal{I},\kappa^{*}} = true \ \text{for every} \ x\text{-variant} \ \kappa^{*} \ \text{of} \ \kappa \\
                  false & : \text{otherwise}
              \end{array}
              \right\}$

    \item	$|\exists x A|^{\mathcal{I},\kappa} = \left\{
              \begin{array}{lr}
                  true  & : |A|^{\mathcal{I},\kappa^{*}} = true \ \text{for some} \ x\text{-variant} \ \kappa^{*} \ \text{of} \ \kappa \\
                  false & : \text{otherwise}
              \end{array}
              \right\}$

\end{enumerate}

\noindent \textbf{Satisfiability of a first-order formula}: A first-order formula $A$ is satisfiable if there is
an interpretation $\mathcal{I}$ and a variable evaluation $\kappa$ for which $|A|^{\mathcal{I},\kappa} = true$
(in which case we say that the interpretation $\mathcal{I}$ and variable evaluation $\kappa$ satisfy $A$),
otherwise unsatisfiable.

If the formula $A$ is closed, its truth value is determined solely by the interpretation. If $|A|^{\mathcal{I}} = true$,
we say that $\mathcal{I}$ satisfies $A$ or otherwise: $\mathcal{I}$ is a model of $A$ ($\mathcal{I} \models A$).\\

\noindent \textbf{Logically true first-order formula}: A first-order logical formula $A$ is logically true
if in every interpretation $\mathcal{I}$ and under every variable evaluation $\kappa$ of $\mathcal{I}$,
$|A|^{\mathcal{I},\kappa} = true$. Its notation: $\models A$.\\

\noindent \textbf{Semantic consequence}: We say that the formula $G$ is a \textit{semantic consequence} of the
formula set $\mathcal{F}$ if for every interpretation $\mathcal{I}$ for which $\mathcal{I} \models \mathcal{F}$
holds, $\mathcal{I} \models G$ is also true (notation: $\mathcal{F} \models G$).\\

\noindent \textbf{Theorem}: Let $A_{1}, A_{2}, ..., A_{n}, B$ ($n \geq 1$) be arbitrary formulas from the same first-order logical language.
Then $\left\{A_{1}, A_{2}, ..., A_{n}\right\} \models B$ if and only if
$A_{1} \wedge A_{2} \wedge ... \wedge A_{n} \wedge \neg B$ is unsatisfiable.\\

\noindent \textbf{Resolution}: Resolution can also be carried out in first-order predicate calculus, and moreover the method
is sound and complete. Difficulty can be caused by the construction of the clauses, which consist of the conjunction of closed,
universally quantified literals. For this our tools are the
prenex and Skolem forms.

\end{document}
